\documentclass[11pt, a4paper]{article}

\usepackage{fontspec}
\setmainfont{Helvetica Neue}
\setmonofont{Menlo}
\usepackage{unicode-math}
\setmathfont{STIX Two Math}

\usepackage[margin=2cm, top=2cm, bottom=2.5cm]{geometry}
\usepackage{microtype}
\usepackage{parskip}
\setlength{\parskip}{0.6em}
\setlength{\parindent}{0pt}

\usepackage[colorlinks=true, linkcolor=NavyBlue, urlcolor=NavyBlue, citecolor=NavyBlue]{hyperref}
\usepackage{xcolor}
\usepackage[dvipsnames]{xcolor}

\usepackage{amsmath}
\usepackage{booktabs}
\usepackage{array}
\usepackage{tabularx}
\newcommand{\thead}[1]{\textbf{#1}}

% --- Graphics ---
\usepackage{graphicx}
\graphicspath{{figures/week30/}}

\usepackage{enumitem}
\setlist[itemize]{leftmargin=1.5em, itemsep=0.2em, topsep=0.3em}
\setlist[enumerate]{leftmargin=1.5em, itemsep=0.2em, topsep=0.3em}

\usepackage[most]{tcolorbox}
\tcbuselibrary{breakable, skins, theorems}

\definecolor{defBg}{HTML}{EAF2FB}
\definecolor{defBd}{HTML}{1F4E79}
\definecolor{exBg}{HTML}{F4F4F4}
\definecolor{exBd}{HTML}{555555}
\definecolor{tipBg}{HTML}{E8F5E9}
\definecolor{tipBd}{HTML}{2E7D32}
\definecolor{trapBg}{HTML}{FCE4EC}
\definecolor{trapBd}{HTML}{AD1457}
\definecolor{pitBg}{HTML}{FFF8E1}
\definecolor{pitBd}{HTML}{B27600}
\definecolor{noteBg}{HTML}{F3E5F5}
\definecolor{noteBd}{HTML}{6A1B9A}
\definecolor{tldrBg}{HTML}{E1F5FE}
\definecolor{tldrBd}{HTML}{0277BD}

\newtcolorbox{defbox}[1][]{enhanced, breakable, colback=defBg, colframe=defBd, arc=2pt, boxrule=0.6pt, left=8pt, right=8pt, top=6pt, bottom=6pt, fonttitle=\bfseries, title={Definition: #1}}
\newtcolorbox{exbox}[1][]{enhanced, breakable, colback=exBg, colframe=exBd, arc=2pt, boxrule=0.6pt, left=8pt, right=8pt, top=6pt, bottom=6pt, fonttitle=\bfseries, title={Worked example: #1}}
\newtcolorbox{exambox}[1][]{enhanced, breakable, colback=tipBg, colframe=tipBd, arc=2pt, boxrule=0.6pt, left=8pt, right=8pt, top=6pt, bottom=6pt, fonttitle=\bfseries, title={Exam-mode answer #1}}
\newtcolorbox{trapbox}[1][]{enhanced, breakable, colback=trapBg, colframe=trapBd, arc=2pt, boxrule=0.6pt, left=8pt, right=8pt, top=6pt, bottom=6pt, fonttitle=\bfseries, title={MCQ trap #1}}
\newtcolorbox{pitfallbox}[1][]{enhanced, breakable, colback=pitBg, colframe=pitBd, arc=2pt, boxrule=0.6pt, left=8pt, right=8pt, top=6pt, bottom=6pt, fonttitle=\bfseries, title={Pitfall #1}}
\newtcolorbox{notebox}[1][]{enhanced, breakable, colback=noteBg, colframe=noteBd, arc=2pt, boxrule=0.6pt, left=8pt, right=8pt, top=6pt, bottom=6pt, fonttitle=\bfseries, title={Lecturer aside #1}}
\newtcolorbox{tldrbox}[1][]{enhanced, breakable, colback=tldrBg, colframe=tldrBd, arc=2pt, boxrule=0.6pt, left=8pt, right=8pt, top=6pt, bottom=6pt, fonttitle=\bfseries, title={TL;DR #1}}

\usepackage{titlesec}
\titleformat{\section}[block]{\Large\bfseries\color{defBd}}{\thesection.}{0.6em}{}
\titleformat{\subsection}[block]{\large\bfseries\color{defBd!80!black}}{\thesubsection}{0.5em}{}
\titleformat{\subsubsection}[block]{\normalsize\bfseries}{}{0pt}{}
\titlespacing*{\section}{0pt}{1.6em}{0.6em}
\titlespacing*{\subsection}{0pt}{1.2em}{0.4em}

\usepackage{fancyhdr}
\pagestyle{fancy}
\fancyhf{}
\fancyhead[L]{\small CM52078 — Week 10}
\fancyhead[R]{\small Probabilistic Reasoning / Bayesian Networks}
\fancyfoot[C]{\small\thepage}
\renewcommand{\headrulewidth}{0.3pt}

\newcommand{\examflag}{\textcolor{tipBd}{\textbf{[EXAM]}}\,}
\newcommand{\trapflag}{\textcolor{trapBd}{\textbf{[TRAP]}}\,}
\newcommand{\doflag}{\textcolor{defBd}{\textbf{[DO]}}\,}

\title{\vspace{-1cm}{\Huge\bfseries Week 10 — Probabilistic Reasoning} \\[0.4em]{\large Independence, Bayesian networks, random variables, expected value}}
\author{\large CM52078 Classical Artificial Intelligence \\[0.2em]\normalsize University of Bath}
\date{Revision notes}

\begin{document}

\maketitle
\thispagestyle{empty}

\vspace{1em}

\begin{tldrbox}
\begin{itemize}[leftmargin=*]
  \item \textbf{Independence}: $A$ and $B$ are independent iff $P(A \mid B) = P(A)$, equivalently $P(A \cap B) = P(A) \cdot P(B)$. Symmetric.
  \item \textbf{Conditional independence}: $A \perp B \mid C$ iff $P(A \mid B, C) = P(A \mid C)$. Weaker than independence — and far more common in real situations.
  \item \textbf{Bayesian network}: a directed acyclic graph (DAG) where nodes are random variables and arrows encode conditional dependence. Each node has a conditional probability table (CPT) over its parents' values.
  \item \textbf{Joint probability via the BN}: $P(x_1, \dots, x_n) = \prod_{i=1}^n P(x_i \mid \mathrm{parents}(X_i))$. \textbf{Massive saving} over storing the full joint — for a 6-variable network, 12 CPT entries instead of 64.
  \item \textbf{Reading independence off a BN}: arrow = dependence; no path = independence; ``$X \to Y \to Z$ with no other route'' = $Z \perp X \mid Y$.
  \item \textbf{Two formal rules}: (1) each node is conditionally independent of its non-descendants given its parents; (2) each node is conditionally independent of all others given its \textbf{Markov blanket} (parents, children, children's parents).
  \item \textbf{Random variable}: a function from the sample space to the real numbers, $X : \Omega \to \mathbb{R}$. Lets us assign \emph{numerical values} to outcomes (e.g.\ payouts, counts).
  \item \textbf{Standard discrete distributions}: \textbf{Bernoulli} $B(p)$ (one binary trial), \textbf{Binomial} $B(n, p)$ ($n$ binary trials).
  \item \textbf{Expected value}: $E(X) = \sum_x x \cdot P(X = x)$. Mean numerical outcome of repeated trials. Compares which gamble or game has the higher long-run payoff.
\end{itemize}
\end{tldrbox}

% =====================================================================
\section{Independence}

\subsection{From the product rule to independence}

Recall the product rule from last week:
\[
P(A \cap B) = P(B) \cdot P(A \mid B).
\]
The conditional appears because $A$ and $B$ might affect one another. \textbf{If they don't}, the conditional collapses:
\[
P(A \mid B) = P(A) \quad \Longleftrightarrow \quad A, B \text{ are independent}.
\]
Substituting back, independence gives the simple product:
\[
P(A \cap B) = P(A) \cdot P(B).
\]

\begin{defbox}[Independence]
Events $A$ and $B$ are \textbf{independent} iff
\[
P(A \mid B) = P(A) \quad \Longleftrightarrow \quad P(A \cap B) = P(A) \cdot P(B).
\]
The relation is \textbf{symmetric}: $A \perp B \iff B \perp A$.
\end{defbox}

\textbf{Genuine independence is rare} in real-world scenarios. Two medical symptoms aren't independent because a single underlying disease can cause both. Two stock prices aren't independent because they share macroeconomic exposure. The variables we model are usually entangled in some web of dependencies.

\subsection{Conditional independence — the workhorse}

\begin{defbox}[Conditional independence]
$A$ is \textbf{conditionally independent} of $B$ given $C$ if
\[
P(A \mid B, C) = P(A \mid C).
\]
Notation: $A \perp B \mid C$. Once you know $C$, knowing $B$ tells you nothing extra about $A$.
\end{defbox}

Note that ``ordinary'' independence is the special case with no conditioning variable. Conditional independence is strictly weaker.

\begin{exbox}[A conditional independence in everyday life]
Let $A$ = ``I get to work on time'', $B$ = ``I leave home on time'', $C$ = ``I catch my train and it's on time.''\\[0.3em]
$A$ and $B$ are \emph{not} independent: leaving late tends to mean arriving late.\\[0.3em]
But $A \perp B \mid C$: if you've already conditioned on whether I caught the train (say I rushed and made it despite leaving late), then whether I left home on time gives no further information about my arrival time.\\[0.3em]
The train-catching status $C$ is the \textbf{reason} for the dependence between $A$ and $B$. Once you know $C$, the dependence vanishes.
\end{exbox}

\textbf{Symptom example.} Two symptoms (cough, rash) are dependent in the population because both correlate with disease status. But \emph{given} the disease status, they may well be independent — once you fix the cause, the symptoms vary on their own.

\subsubsection*{Equivalent forms of conditional independence}

The following are all equivalent:
\begin{itemize}
  \item $A \perp B \mid C$.
  \item $P(A \mid B, C) = P(A \mid C)$.
  \item $P(B \mid A, C) = P(B \mid C)$.
  \item $P(A, B \mid C) = P(A \mid C) \cdot P(B \mid C)$.
\end{itemize}

The fourth form (factorisation given $C$) is what makes Bayesian networks tractable.

% =====================================================================
\section{Bayesian networks}

A \textbf{Bayesian network} is a graphical language for capturing dependence, independence, and conditional independence between many random variables in one compact picture.

\subsection{The structure}

\begin{defbox}[Bayesian network]
A directed acyclic graph (\textbf{DAG}) plus numeric parameters at each node:
\begin{itemize}[itemsep=0pt, topsep=0pt]
  \item Nodes are random variables.
  \item Directed edges encode \emph{direct} dependence (often interpreted as causation, though formally it's just statistical conditioning).
  \item No directed cycles allowed — you can't follow arrows back to where you started.
  \item Each node has a \textbf{conditional probability table (CPT)} giving $P(X \mid \mathrm{parents}(X))$.
\end{itemize}
\end{defbox}

\textbf{Reading dependencies off the graph (introductory rules):}
\begin{itemize}
  \item \textbf{Arrow between $X$ and $Y$} $\Rightarrow$ $X$ and $Y$ are dependent.
  \item \textbf{No path between $X$ and $Y$} $\Rightarrow$ $X$ and $Y$ are independent.
  \item \textbf{Path $X \to Y \to Z$ with no other route from $X$ to $Z$} $\Rightarrow$ $Z \perp X \mid Y$.
\end{itemize}

(These rules don't catch \emph{all} conditional independences. For full coverage you'd use d-separation; this lecture stops at the simpler rules below.)

\begin{center}

\footnotesize\textit{The work-commute Bayesian network as a DAG. Each ellipse is a binary random variable; each arrow points from cause to effect and encodes \textbf{direct dependence}. Apply the three reading rules here: $A$ and $D$ are joined by an arrow, so they are dependent; $C$ (the moon phase) has no arrow to anything, so it is independent of every other variable; and the only route from $B$ to $F$ runs $B \to D \to F$, so $F \perp B \mid D$ --- once you know whether the train was caught, the roadworks tell you nothing more about arriving on time. Notice the graph is acyclic: you can never follow arrows back to a node you started from.}
\end{center}

\subsection{Three canonical patterns}

Three sub-patterns appear in every BN. Knowing them helps quick reading of conditional independence.

\subsubsection*{Pattern 1: Chain ($X \to Y \to Z$)}

\begin{itemize}
  \item $X$ and $Z$ are dependent (information flows $X \to Y \to Z$).
  \item $X \perp Z \mid Y$. Once $Y$ is known, the chain is broken.
\end{itemize}

\subsubsection*{Pattern 2: Common cause / fork ($X \leftarrow Y \to Z$)}

\begin{itemize}
  \item $X$ and $Z$ are dependent (common cause $Y$ correlates them).
  \item $X \perp Z \mid Y$. Once $Y$ is known, the correlation is explained.
\end{itemize}

\subsubsection*{Pattern 3: Common effect / collider / V-structure ($X \to Y \leftarrow Z$)}

\begin{itemize}
  \item $X$ and $Z$ are independent unconditionally.
  \item $X$ and $Z$ are \emph{dependent} given $Y$. Conditioning on a common effect makes its causes correlated. (``Explaining away'': if I see the effect and one cause, I can rule out the other.)
\end{itemize}

\textbf{This third pattern is the trickiest.} It's the reason the Markov blanket includes children's parents.

\begin{exbox}[Why conditioning on a collider couples its parents]
$X = $ ``it rained,'' $Z = $ ``the sprinkler ran,'' $Y = $ ``the lawn is wet.'' Either independent cause turns the lawn wet; both are usually independent of each other (rain doesn't depend on sprinkler activity).\\[0.3em]
\textbf{Unconditional:} $P(X) = 0.3$, $P(Z) = 0.1$, and $X \perp Z$ — knowing it rained tells you nothing about whether the sprinkler ran.\\[0.3em]
\textbf{Now condition on $Y = T$ (lawn is wet).} The lawn became wet for some reason. If you additionally learn $X = T$ (it rained), that already explains the wetness — so $P(Z \mid Y, X) < P(Z \mid Y)$. The two causes have become \emph{negatively correlated} given the effect, even though they were unconditionally independent. This is ``explaining away.''\\[0.3em]
\textbf{Why the algebra forces this.} Conditional on $Y$, $P(X, Z \mid Y) = P(X) P(Z) P(Y \mid X, Z) / P(Y)$ — and that final factor $P(Y \mid X, Z)$ does not factor as a function of $X$ alone times a function of $Z$ alone. So $X$ and $Z$ cannot be independent given $Y$. \emph{Conditioning on a child of two parents creates a back-door path through the child that re-couples them.}\\[0.3em]
\textbf{Consequence for the Markov blanket.} To make a node $X$ conditionally independent of every other node, you must block this back-door. The blocking set must include each of $X$'s children (so paths through them are visible) \emph{and} every other parent of those children (so the V-structure-induced coupling is severed). That's where the children's-parents term in the blanket comes from.
\end{exbox}

\subsection{Running example}

Six binary variables for whether you arrive at work on time:

\begin{exbox}[The work-commute Bayesian network]
\begin{itemize}[itemsep=0pt, topsep=0pt]
  \item $A$ = I leave my house on time.
  \item $B$ = There are roadworks on the way to the station.
  \item $C$ = It is a waning gibbous moon.
  \item $D$ = I catch my train.
  \item $E$ = My train is delayed.
  \item $F$ = I arrive at work on time.
\end{itemize}
Edges: $A \to D$, $B \to D$, $D \to F$, $E \to F$. Node $C$ is isolated.\\[0.3em]
\textbf{Implied independence:} $C$ is independent of all other variables (it has no path).\\[0.3em]
\textbf{Implied conditional independence:} $F \perp B \mid D$ (the only route from $B$ to $F$ goes through $D$).
\end{exbox}

\subsection{What a Bayesian network represents: the joint probability}

\begin{defbox}[Joint probability factorisation]
A Bayesian network defines a joint distribution over all its variables:
\[
P(x_1, x_2, \dots, x_n) = \prod_{i=1}^n P\bigl(x_i \mid \mathrm{parents}(X_i)\bigr).
\]
\end{defbox}

For our example:
\[
P(A, B, C, D, E, F) = P(A) \cdot P(B) \cdot P(C) \cdot P(D \mid A, B) \cdot P(E) \cdot P(F \mid D, E).
\]

\begin{center}

\footnotesize\textit{The general chain rule for Bayesian networks (highlighted box) alongside the work-commute graph it specialises to. The product $\prod_{i=1}^n p(x_i \mid \mathrm{parents}(X_i))$ says: take one factor per variable, and condition each variable on its parents only. The graph fixes what ``parents'' means --- the nodes with arrows pointing \emph{into} the variable. This single formula is the engine of every Bayes-net calculation: read the parents off the picture, write the matching conditional, multiply.}
\end{center}

\begin{center}

\footnotesize\textit{The same network, now read as a recipe for the joint distribution. Each node contributes exactly one factor to the product, and that factor conditions \textbf{only on the node's parents}. Trace it: $A$, $B$, $C$, $E$ have no incoming arrows, so they contribute plain marginals $P(A)$, $P(B)$, $P(C)$, $P(E)$; $D$ has parents $A$ and $B$, so it contributes $P(D \mid A, B)$; $F$ has parents $D$ and $E$, so it contributes $P(F \mid D, E)$. Multiply the six factors and you have the full joint $P(A,B,C,D,E,F)$ --- this is the chain rule for Bayes nets, and the conditional independences baked into the graph are what let the conditioning sets stay so small.}
\end{center}

\textbf{Read this carefully.} Variables with no parents ($A, B, C, E$) need only marginal probabilities. Variables with parents need conditionals over those parent values only — \emph{not} over every other variable in the network.

\subsection{Why the BN factorisation saves so much storage}

For 6 binary variables, the full joint has $2^6 = 64$ entries (each row a different combination of T/F). \textbf{Storage is exponential in the number of variables.}

\begin{center}

\footnotesize\textit{The full joint probability table $P(A,B,C,D,E,F)$ stored naively, with no use of network structure. Each of the six binary variables doubles the row count, so the table has $2^6 = 64$ rows --- one per distinct True/False assignment, from all-False at the top to all-True at the bottom. The rightmost column holds one probability value per row. This is the table the BN factorisation lets us avoid: it grows exponentially, so for a few dozen variables it is already too large to store or fill in by hand.}
\end{center}

Using the BN factorisation:
\begin{itemize}
  \item $P(A=T)$: 1 entry (with $P(A=F) = 1 - P(A=T)$).
  \item $P(B=T)$: 1 entry. $P(C=T)$: 1. $P(E=T)$: 1.
  \item $P(D=T \mid A, B)$: 4 entries (one per combination of $A, B$).
  \item $P(F=T \mid D, E)$: 4 entries.
\end{itemize}
Total: $1 + 1 + 1 + 1 + 4 + 4 = 12$ entries. \textbf{12 vs.\ 64 — a factor-5 saving for a tiny network.}

\begin{center}

\footnotesize\textit{The same joint distribution stored as the network's six conditional probability tables (CPTs) instead of one 64-row table. The four parentless nodes $A$, $B$, $C$, $E$ each need a single number, $P(\cdot = \text{True})$ --- the False probability is fixed by ``probabilities sum to 1.'' The two nodes with parents are larger: $P(D = \text{True} \mid A, B)$ has one row per combination of its two parents ($2^2 = 4$ rows) and likewise $P(F = \text{True} \mid D, E)$ has 4. Counting only the highlighted probability cells gives $1+1+1+1+4+4 = 12$ values --- the ``12 elements'' the slide flags, versus 64 for the naive table.}
\end{center}

For real networks with hundreds of variables, the saving is the difference between feasible and infeasible. \textbf{This compactness is the whole point of Bayesian networks.}

\subsubsection*{General formula for CPT entries}

For each binary node $X$ with $|\mathrm{parents}(X)| = k$:
\[
\text{CPT entries needed} = 2^k.
\]

(Why $2^k$? For each combination of parent values — $2^k$ of them — we need the conditional probability $P(X = T \mid \text{parents})$. The $X = F$ case is determined by ``probabilities sum to 1.'')

For a $k$-valued node with binary parents: $2^{|\mathrm{parents}|}$ parent combinations $\times$ $(k - 1)$ entries each.

\examflag{}Past papers like to ask you to count CPT entries given a network and binary variables. The trick: each variable contributes $2^{|\mathrm{parents}|}$ entries.

\subsection{Worked numbers: the alarm network}

The lecturer's classic textbook example (Russell \& Norvig). Five binary variables: \emph{Burglary}, \emph{Earthquake}, \emph{Alarm}, \emph{JohnCalls}, \emph{MaryCalls}. Edges: Burglary $\to$ Alarm, Earthquake $\to$ Alarm, Alarm $\to$ JohnCalls, Alarm $\to$ MaryCalls.

CPTs (truncated):
\[
P(B = T) = 0.001, \quad P(E = T) = 0.002.
\]
\[
P(A = T \mid B, E):\ TT = 0.95,\ TF = 0.94,\ FT = 0.29,\ FF = 0.001.
\]
\[
P(J = T \mid A): A = T \to 0.90, A = F \to 0.05.
\]
\[
P(M = T \mid A): A = T \to 0.70, A = F \to 0.01.
\]

\begin{center}

\footnotesize\textit{The classic alarm Bayesian network with its CPTs attached. \emph{Burglary} and \emph{Earthquake} are parentless, so each has a one-row table giving $P(\cdot = \text{true})$. \emph{Alarm} has both as parents: its CPT has a row for every $B,E$ combination ($tt, tf, ft, ff$), e.g.\ the alarm fires with probability $0.95$ when both occur but only $0.001$ on its own. \emph{JohnCalls} and \emph{MaryCalls} each depend only on \emph{Alarm}, so their tables have two rows ($A = t$ or $f$). The worked sum below multiplies one selected entry from each table; this is exactly the chain-rule factorisation $P(j,m,a,\neg b,\neg e) = P(j \mid a)P(m \mid a)P(a \mid \neg b, \neg e)P(\neg b)P(\neg e)$ read straight off the graph. Image credit: Russell \& Norvig, AIMA.}
\end{center}

\begin{exbox}[Computing $P(j, m, a, \neg b, \neg e)$]
The joint probability of John and Mary both calling, alarm going off, but no burglary and no earthquake:
\begin{align*}
P(j, m, a, \neg b, \neg e) &= P(j \mid a) \cdot P(m \mid a) \cdot P(a \mid \neg b, \neg e) \cdot P(\neg b) \cdot P(\neg e) \\
&= 0.90 \cdot 0.70 \cdot 0.001 \cdot 0.999 \cdot 0.998 \\
&\approx 0.000628.
\end{align*}
\textbf{The probability of a false alarm scenario where both neighbours call: about 0.06\%.}
\end{exbox}

\subsubsection*{CPT entry count for the alarm network}

\begin{itemize}
  \item $B$: no parents. $2^0 = 1$ entry.
  \item $E$: no parents. $1$ entry.
  \item $A$: 2 parents ($B$, $E$). $2^2 = 4$ entries.
  \item $J$: 1 parent ($A$). $2^1 = 2$ entries.
  \item $M$: 1 parent ($A$). $2$ entries.
\end{itemize}
\textbf{Total: $1 + 1 + 4 + 2 + 2 = 10$ entries.} The full joint over 5 binary variables would have $2^5 = 32$. Saving: a factor of about 3.

% =====================================================================
\section{Reading conditional independences off a BN}

The introductory rules don't handle every case. Two stronger rules cover most past-paper scenarios.

\begin{center}

\footnotesize\textit{The two formal conditional-independence rules, side by side, for the same target node $X$ (dark grey). \textbf{Left --- non-descendants rule:} shading the $U$ nodes ($X$'s parents) makes $X$ independent of all its non-descendants; the $Z$ nodes are non-descendants and the shaded parents block every path to them. \textbf{Right --- Markov blanket:} the grey ring is $X$'s parents ($U$), its children ($Y$), \emph{and} its children's other parents ($Z$); conditioning on this whole ring makes $X$ independent of everything outside it. Note why $Z$ (the children's parents) must be in the ring: $X \to Y \leftarrow Z$ is a collider, so conditioning on the child $Y$ couples $X$ to $Z$, and that new path must be blocked. Image credit: Russell \& Norvig, AIMA.}
\end{center}

\subsection{Rule 1: non-descendants given parents}

\begin{defbox}[Non-descendants rule]
Each node is conditionally independent of its \textbf{non-descendants} given its \textbf{parents}.\\[0.3em]
A non-descendant of $X$ is any node not reachable by following arrows forward from $X$.
\end{defbox}

In the work-commute example: $D$'s parents are $A$ and $B$. $D$'s non-descendants are everything except $F$ (i.e., $A$, $B$, $C$, $E$). Given $\{A, B\}$, $D$ is independent of $C$ and $E$.

\textbf{Why does this hold?} Because in the BN's joint factorisation, $D$'s factor depends only on $A$ and $B$ (its parents). Conditioning on $A, B$ fixes everything that affects $D$'s distribution.

\subsection{Rule 2: the Markov blanket}

\begin{defbox}[Markov blanket]
The \textbf{Markov blanket} of node $X$ consists of:
\begin{itemize}[itemsep=0pt, topsep=0pt]
  \item $X$'s parents,
  \item $X$'s children,
  \item $X$'s children's parents (i.e., the spouses of $X$'s children — the other parents of any of $X$'s children).
\end{itemize}
\textbf{Each node is conditionally independent of every other node given its Markov blanket.}
\end{defbox}

\textbf{The asymmetry that catches everyone:} the blanket includes children's \emph{parents}, but not parents' \emph{children}. (Why? Because parents already block paths upward, but two paths converging at a common child create a new dependency you have to block — that's the V-structure / collider pattern from earlier.)

\begin{exbox}[Markov blanket in the work-commute example]
For node $D$:
\begin{itemize}[itemsep=0pt, topsep=0pt]
  \item Parents: $A$, $B$.
  \item Children: $F$.
  \item Children's parents (other than $D$): $E$ (the other parent of $F$).
\end{itemize}
$D$'s Markov blanket = $\{A, B, F, E\}$. Given this set, $D$ is independent of any other node in the network — in this case, just $C$.\\[0.3em]
For node $E$: parents = $\emptyset$, children = $\{F\}$, children's parents = $\{D\}$. Blanket = $\{F, D\}$. Given $D$ and $F$, $E$ is independent of $A$, $B$, $C$.
\end{exbox}

\subsubsection*{Worked example: Markov blanket of $A$ in the alarm network}

\begin{exbox}[Markov blanket of $A$ (Alarm)]
\textbf{Network:} $B \to A$, $E \to A$, $A \to J$, $A \to M$.\\[0.3em]
\begin{itemize}[itemsep=0pt, topsep=0pt]
  \item Parents of $A$: $B$, $E$.
  \item Children of $A$: $J$, $M$.
  \item Children's parents (other than $A$): none — $J$ has only $A$ as parent, $M$ has only $A$.
\end{itemize}
\textbf{Markov blanket of $A$ = $\{B, E, J, M\}$.} Given these four, $A$ is independent of \dots\ well, there are no other variables in this 5-node network. So in this small example, the blanket is the whole rest of the graph.
\end{exbox}

\begin{trapbox}[: ``children's parents'' is asymmetric]
A favourite past-paper trap. The blanket includes \emph{children's parents}, NOT \emph{parents' children}. Don't get tripped up — try writing the blanket on a small example before answering.\\[0.3em]
\textbf{Why the asymmetry exists.} It's the V-structure / collider pattern: $X \to Y \leftarrow Z$. Without conditioning on $Y$, $X$ and $Z$ are independent. Conditioning on $Y$ creates dependence between them. So to fully isolate a node from the rest of the graph, you need to block these new dependencies — which means including the spouses of your children in the blanket.
\end{trapbox}

% =====================================================================
\section{Random variables}

So far we've talked about events and outcomes abstractly. \textbf{Random variables} promote outcomes to numerical values, which lets us compute averages and use standard distributions.

\subsection{Definition}

\begin{defbox}[Random variable]
A \textbf{random variable} is a function from the sample space to the real numbers, $X : \Omega \to \mathbb{R}$. It assigns a numerical value to each outcome.
\end{defbox}

\textbf{Important.} The same probabilistic experiment can have many random variables on it.

\begin{exbox}[Two random variables on a two-coin-toss experiment]
\textbf{Sample space:} $\{HH, HT, TH, TT\}$, fair coin.\\[0.3em]
\textbf{Random variable $X$:} ``count the heads.'' $X(HH) = 2, X(HT) = 1, X(TH) = 1, X(TT) = 0$.\\[0.3em]
\textbf{Random variable $W$:} ``gambling game where you start with $\pounds 1$ and each heads doubles, each tails halves.'' $W(HH) = 4, W(HT) = 1, W(TH) = 1, W(TT) = \frac{1}{4}$.\\[0.3em]
\textbf{Same experiment, different random variables, different expected values.} The choice of random variable depends on what you want to measure.
\end{exbox}

\subsection{Discrete vs.\ continuous}

\begin{defbox}[Discrete vs.\ continuous random variable]
A random variable is \textbf{discrete} if its range is finite or countably infinite (e.g., dice rolls, integer counts). It is \textbf{continuous} if its range is uncountably infinite (e.g., heights, weights, real-valued measurements).
\end{defbox}

The mathematics of discrete and continuous random variables is different (continuous needs probability density functions and integrals). \textbf{This unit focuses on discrete} — the random variables in Bayesian networks (binary truth values) are discrete.

\subsection{Standard discrete distributions}

\subsubsection*{Bernoulli distribution $B(p)$}

A single binary trial with probability $p$ of ``success.''
\[
P(X = 1) = p, \qquad P(X = 0) = 1 - p.
\]
$X \sim B(p)$.

Coin flip with probability $p$ of heads. Each binary node in a Bayes net (conditioning on parents) has Bernoulli-distributed values.

\subsubsection*{Binomial distribution $B(n, p)$}

$n$ independent Bernoulli trials with success probability $p$. The random variable counts the number of successes:
\[
P(X = k) = \binom{n}{k} p^k (1 - p)^{n - k}.
\]
$X \sim B(n, p)$.

For $n = 3$:
\[
P(X = 0) = (1-p)^3, \quad P(X = 1) = 3 p (1-p)^2, \quad P(X = 2) = 3 p^2 (1-p), \quad P(X = 3) = p^3.
\]

\subsubsection*{Other distributions to know in passing}

\begin{itemize}
  \item \textbf{Categorical}: a single roll of a $k$-sided die. Each outcome has its own probability $p_i$ with $\sum p_i = 1$.
  \item \textbf{Multinomial}: $n$ rolls of a $k$-sided die. Generalises binomial.
  \item \textbf{Poisson}: count of rare events in a fixed interval. Rate $\lambda$. $P(X = k) = e^{-\lambda} \lambda^k / k!$.
  \item \textbf{Geometric}: number of Bernoulli trials until the first success. $P(X = k) = (1-p)^{k-1} p$.
\end{itemize}

\textbf{Out of scope for this exam} but worth recognising in passing.

% =====================================================================
\section{Expected value}

\subsection{Definition and rationale}

\begin{defbox}[Expected value (expectation)]
For a discrete random variable $X$:
\[
E(X) = \sum_x x \cdot P(X = x).
\]
The long-run average value of $X$ if the experiment is repeated many times.
\end{defbox}

\textbf{Why we care.} Random variables let us assign numerical values to outcomes; expectation gives a single summary number for ``what to expect on average.'' It's how we compare gambles, choose between policies, and value uncertain payoffs.

\textbf{Linearity of expectation.} A useful property: $E(aX + bY) = a E(X) + b E(Y)$, regardless of whether $X$ and $Y$ are independent.

\subsection{Two simple examples}

\begin{exbox}[Coin-flip and dice-roll expectations]
\textbf{Coin flip} with $X(H) = 1, X(T) = 0$. Fair coin:
\[
E(X) = 1 \cdot \tfrac{1}{2} + 0 \cdot \tfrac{1}{2} = 0.5.
\]
\textbf{Dice roll} with $X(\text{face}) = \text{face number}$:
\[
E(X) = \sum_{x = 1}^{6} x \cdot \tfrac{1}{6} = \frac{1 + 2 + 3 + 4 + 5 + 6}{6} = \frac{21}{6} = 3.5.
\]
The expected value of a single die roll is 3.5. \textbf{Note: 3.5 isn't a possible single outcome.} Expectation is an average over many trials.
\end{exbox}

\subsection{Expectation depends on the random variable, not just the experiment}

\begin{exbox}[Comparing two gambles on the same experiment]
Two-coin-toss experiment, fair coin. Compare:\\[0.3em]
\textbf{Variable $X$ (count heads):} $X(HH) = 2, X(HT) = X(TH) = 1, X(TT) = 0$. Each outcome with probability $\frac{1}{4}$.
\[
E(X) = \frac{2 + 1 + 1 + 0}{4} = 1.
\]
\textbf{Variable $W$ (start with $\pounds 1$, heads doubles, tails halves):} $W(HH) = 4, W(HT) = W(TH) = 1, W(TT) = \frac{1}{4}$.
\[
E(W) = \frac{4 + 1 + 1 + 0.25}{4} = \frac{6.25}{4} = \frac{25}{16} \approx 1.56.
\]
\textbf{Game $W$ has the higher expected value} — you'd rather play it. Despite the same sample space and probabilities, the random variable changes the answer.
\end{exbox}

\textbf{Standard MCQ pattern:} ``which game would you prefer?'' or ``what is the expected fine?'' — compute the expectation by summing $x \cdot P(X = x)$ over all values.

\begin{exbox}[Expected fine for parking in an alley]
You park your scooter; there's a 5\% chance an inspector walks past and gives you a $\pounds 60$ fine. What's the expected fine?
\[
E(\text{Fine}) = 60 \cdot 0.05 + 0 \cdot 0.95 = 3.
\]
$\pounds 3$ — long-run average cost per parking. (Decision rule: park here only if convenience is worth more than $\pounds 3$.)
\end{exbox}

\subsubsection*{Expected reward from a probability table}

\begin{exbox}[{Expected reward from a joint probability table}]
Joint $P(\text{Reward}, \text{State})$:\\
\begin{center}
\begin{tabular}{@{}lcc@{}}
\toprule
 & State 1 & State 2 \\
\midrule
Reward = 5  & 0.1 & 0.3 \\
Reward = 10 & 0.5 & 0.1 \\
\bottomrule
\end{tabular}
\end{center}
\textbf{First, marginalise to get $P(\text{Reward})$:}
\[
P(\text{Reward} = 5) = 0.1 + 0.3 = 0.4, \quad P(\text{Reward} = 10) = 0.5 + 0.1 = 0.6.
\]
\textbf{Then expected reward:}
\[
E(\text{Reward}) = 5 \cdot 0.4 + 10 \cdot 0.6 = 2 + 6 = 8.
\]
\end{exbox}

% =====================================================================
\section{Naive Bayes}

A particularly useful Bayes-net topology shows up so often it has its own name.

\begin{defbox}[Naive Bayes classifier]
A Bayesian network with one ``class'' node $C$ and many ``feature'' nodes $F_1, F_2, \dots, F_n$, all of which are children of $C$. The features are assumed \emph{conditionally independent given the class}.\\[0.3em]
The factorised joint:
\[
P(C, F_1, \dots, F_n) = P(C) \cdot \prod_{i=1}^n P(F_i \mid C).
\]
\end{defbox}

\textbf{Why ``naive''?} Because the conditional-independence assumption is usually false in practice — features are usually correlated even given the class. But the assumption keeps the model tractable, and naive Bayes often performs surprisingly well as a baseline classifier.

\subsubsection*{Use in classification}

To classify a new instance with features $f_1, \dots, f_n$, compute:
\[
P(C = c \mid f_1, \dots, f_n) \propto P(C = c) \cdot \prod_{i=1}^n P(f_i \mid C = c).
\]
Pick the class $c$ that maximises this. Common in spam filters, medical diagnosis, sentiment analysis.

\subsubsection*{Worked example: spam filtering}

\begin{exbox}[Naive Bayes spam classification]
Class $S \in \{\text{spam}, \text{ham}\}$. Three features:
\begin{itemize}[itemsep=0pt, topsep=0pt]
  \item $K$: contains keyword from watchlist.
  \item $R$: positive sender reputation.
  \item $L$: long message.
\end{itemize}
\textbf{Joint via naive Bayes:}
\[
P(K, R, L, S) = P(S) \cdot P(K \mid S) \cdot P(R \mid S) \cdot P(L \mid S).
\]
Given a message with $K = 1, R = 0, L = 0$, classify by computing both $P(\text{spam} \mid K, R, L)$ and $P(\text{ham} \mid K, R, L)$ and picking the larger.
\end{exbox}

% =====================================================================
\section{Inference in Bayes nets (briefly)}

The whole point of building a BN is to answer queries like ``given evidence $E$, what's the probability of variable $X$ being true?''

\begin{itemize}
  \item \textbf{Exact inference} via variable elimination or junction tree algorithms. Works for small networks; complexity $O(2^n)$ in the worst case (NP-hard in general).
  \item \textbf{Approximate inference} via sampling (Gibbs sampling, MCMC, particle filtering). Works for large networks at the cost of accuracy.
\end{itemize}

\textbf{Out of scope for this exam.} For exam purposes you need to:
\begin{itemize}
  \item Read the joint factorisation off the network.
  \item Compute joint probabilities of specific configurations using CPTs.
  \item Apply Bayes' theorem manually for small queries.
  \item Identify conditional independences via the rules above.
\end{itemize}

% =====================================================================
\section{Exam-mode answers}

\begin{exambox}[(MCQ): Which of the following Bayesian networks describes the situation: ``Weather causes traffic and umbrella; traffic causes delay; nothing causes weather''?]
\textit{Variables: $W$ (Weather), $T$ (Traffic), $U$ (Umbrella), $D$ (Delay).}\\[0.3em]
\textit{Correct edges: $W \to T$, $W \to U$, $T \to D$.}\\[0.3em]
\textit{Trap distractors usually have edges in wrong direction or extra/missing edges.}
\end{exambox}

\begin{exambox}[(MCQ): For the BN above, which of the following is the correct factorised joint probability?]
\textit{Apply $P(x_1, \dots, x_n) = \prod P(x_i \mid \mathrm{parents}(X_i))$.}
\[
P(W, T, U, D) = P(W) \cdot P(T \mid W) \cdot P(U \mid W) \cdot P(D \mid T).
\]
\textit{Trap: $P(W, T, U, D) = P(W) P(T) P(U) P(D)$ — assumes everything's independent (false). $P(W \mid T, U, D) P(T \mid U, D) P(U \mid D) P(D)$ — generic chain rule, doesn't exploit conditional independences from the network.}
\end{exambox}

\begin{exambox}[(MCQ): Consider a Bayes network with all variables binary. The minimum total number of elements required to store the joint distribution, using the BN factorisation? (Network: $W \to T$, $W \to U$, $T \to D$, $D$ has no other parents, $W$ has no parents.)]
\textit{For each binary node, store $2^{|\mathrm{parents}|}$ entries (one per parent combination, given the other entry is $1 - p$):}
\begin{itemize}[leftmargin=*, itemsep=0pt, topsep=0pt]
  \item $W$ has no parents: $2^0 = 1$ entry.
  \item $T$ has 1 parent ($W$): $2^1 = 2$ entries.
  \item $U$ has 1 parent ($W$): $2$ entries.
  \item $D$ has 1 parent ($T$): $2$ entries.
\end{itemize}
\textit{Total: $1 + 2 + 2 + 2 = 7$ entries.}
\end{exambox}

\begin{exambox}[(MCQ): Which of the following statements about a Bayesian network are true? Mark all that apply.]
\textit{TRUE:}
\begin{itemize}[leftmargin=*, itemsep=0pt, topsep=0pt]
  \item ``A node is conditionally independent of its non-descendants given its parents.''
  \item ``Each node is conditionally independent of all others given its Markov blanket.''
  \item ``A Bayesian network can compactly represent a joint distribution that would otherwise require an exponentially large table.''
\end{itemize}
\textit{FALSE: ``Two nodes connected by a directed edge are always independent'' (false — they're directly dependent).}
\end{exambox}

\begin{exambox}[(MCQ): A patient presents with three symptoms (fever, cough, fatigue). Assume conditional independence of the symptoms given flu. Which expression gives $P(\mathrm{flu} \mid \mathrm{fever}, \mathrm{cough}, \mathrm{fatigue})$?]
\textit{Bayes plus conditional independence (Naive Bayes):}
\[
P(\mathrm{flu} \mid \mathrm{symptoms}) = \frac{P(\mathrm{fever} \mid \mathrm{flu}) P(\mathrm{cough} \mid \mathrm{flu}) P(\mathrm{fatigue} \mid \mathrm{flu}) P(\mathrm{flu})}{P(\mathrm{fever}, \mathrm{cough}, \mathrm{fatigue})}.
\]
\textit{Conditional independence given flu lets you factor the joint $P(\text{symptoms} \mid \text{flu})$ into individual products. This is the Naive Bayes form.}
\end{exambox}

% =====================================================================
\section*{Key equations / algorithms cheat sheet}

\textbf{Independence:}
\[
A \perp B \iff P(A \mid B) = P(A) \iff P(A \cap B) = P(A) \cdot P(B).
\]

\textbf{Conditional independence:}
\[
A \perp B \mid C \iff P(A \mid B, C) = P(A \mid C).
\]

\textbf{Bayesian network factorisation:}
\[
P(x_1, \dots, x_n) = \prod_{i=1}^n P\bigl(x_i \mid \mathrm{parents}(X_i)\bigr).
\]

\textbf{CPT entries per binary node:} $2^{|\mathrm{parents}|}$ (storing $P(X = T \mid \cdot)$; the $F$ entry is $1 -$ that).

\textbf{Conditional independence rules:}
\begin{itemize}[itemsep=0pt, topsep=0pt]
  \item Each node $X$ $\perp$ its non-descendants $\mid$ its parents.
  \item Each node $X$ $\perp$ all other nodes $\mid$ its Markov blanket = (parents, children, children's parents).
\end{itemize}

\textbf{Three sub-patterns:}
\begin{itemize}[itemsep=0pt, topsep=0pt]
  \item Chain $X \to Y \to Z$: $X \perp Z \mid Y$.
  \item Fork $X \leftarrow Y \to Z$: $X \perp Z \mid Y$.
  \item Collider $X \to Y \leftarrow Z$: $X \perp Z$ unconditionally; $X \not\perp Z \mid Y$.
\end{itemize}

\textbf{Random variable:} $X : \Omega \to \mathbb{R}$.

\textbf{Standard distributions:}
\[
\text{Bernoulli } B(p): P(X = 1) = p, \; P(X = 0) = 1 - p.
\]
\[
\text{Binomial } B(n, p): P(X = k) = \binom{n}{k} p^k (1 - p)^{n - k}.
\]

\textbf{Expected value (discrete):}
\[
E(X) = \sum_x x \cdot P(X = x).
\]

\textbf{Linearity of expectation:}
\[
E(aX + bY) = a E(X) + b E(Y) \text{ regardless of independence}.
\]

\textbf{Naive Bayes:}
\[
P(C, F_1, \dots, F_n) = P(C) \prod_i P(F_i \mid C).
\]

% =====================================================================
\section*{Pitfalls and MCQ traps consolidated}

\begin{trapbox}[summary]
\begin{tabular}{@{}p{6.5cm}p{8cm}@{}}
\toprule
\thead{Trap} & \thead{Right answer} \\
\midrule
``Independent $\Leftrightarrow$ mutually exclusive'' & \textbf{Opposite!} ME means $A \cap B = \emptyset$. Independence is $P(A \cap B) = P(A)P(B)$. \\
``$P(A \cap B) = P(A) P(B)$ always''               & \textbf{Only if independent.} \\
``All variables in a real-world model are independent'' & \textbf{Rarely.} Conditional independence is more common. \\
Markov blanket = parents + children                  & \textbf{Plus children's parents.} \\
Bayes net edges represent causation                   & Often, but \textbf{strictly they encode dependence}. Reverse direction is sometimes equivalent. \\
``Joint = $P(A) P(B) \dots$ when network has structure'' & \textbf{No}, multiply each $P(X \mid \mathrm{parents}(X))$. \\
Random variables are events                            & \textbf{No}, they are functions $\Omega \to \mathbb{R}$. \\
Expected value is the most-likely outcome              & \textbf{No}, it's the long-run average. \\
``$X \to Y \leftarrow Z$ implies $X \perp Z$ given $Y$'' & \textbf{No, the opposite!} Collider conditioning creates dependence. \\
\bottomrule
\end{tabular}
\end{trapbox}

\textbf{Other confusions:}

\begin{itemize}
  \item \textbf{Bayes net direction.} The lecturer used ``causation'' as a teaching analogy. Strictly, the network just encodes which variables are conditionally dependent. Two networks with reversed edges \emph{can} represent the same joint, but with different parametrisations.
  \item \textbf{$2^{|\mathrm{parents}|}$ vs.\ $2^{|\mathrm{parents}|} \cdot 2$.} Past papers ask for ``minimum entries.'' For binary nodes, $2^{|\mathrm{parents}|}$ entries suffice (the other half is determined by ``probabilities sum to 1''). For $k$-valued nodes, it's $2^{|\mathrm{parents}|} \cdot (k - 1)$.
  \item \textbf{Discrete vs.\ continuous random variables.} The Bayesian-network examples in this lecture all use discrete (typically binary) variables. Continuous BNs are a separate topic.
  \item \textbf{Naive Bayes.} A common Bayes-net topology in the past papers: a ``class'' node points to many ``feature'' nodes, with the features assumed conditionally independent given the class. The full joint is $P(\text{class}) \prod P(\text{feature}_i \mid \text{class})$.
  \item \textbf{Expected value is linear.} $E(aX + bY) = a E(X) + b E(Y)$ — useful for combining payoffs or losses across multiple variables. \emph{Doesn't require independence.}
  \item \textbf{Don't confuse expectation with mode.} The expectation $E(X) = 3.5$ for a die isn't a possible single roll; it's the average across many rolls.
  \item \textbf{Three sub-patterns to memorise.} Chain, fork, collider. The first two have the same independence structure ($X \perp Z \mid Y$); the collider has the opposite ($X \perp Z$ \emph{without} conditioning, $X \not\perp Z$ \emph{with}).
\end{itemize}

% =====================================================================
\section*{Self-check questions}

\begin{enumerate}
  \item \textbf{Distinguish} independence and mutual exclusivity. Give an example of each.
  \item \textbf{Define} conditional independence. Construct a small example where $A$ and $B$ are dependent, but $A \perp B \mid C$.
  \item \textbf{Explain} why genuine independence between two variables is rare in real-world settings, but conditional independence is common.
  \item \textbf{Describe} the three sub-patterns (chain, fork, collider) and the conditional independence each implies (or breaks).
  \item \textbf{Describe} a Bayesian network. What information is encoded in the graph, and what is encoded in the conditional probability tables?
  \item \textbf{Compute} the number of entries needed to store the full joint over 6 binary variables. Compare to the BN factorisation for a network with the parent-counts $\{0, 0, 0, 2, 0, 2\}$.
  \item \textbf{Apply} the joint-factorisation formula to a 4-node BN with edges $W \to T$, $W \to U$, $T \to D$. Write out $P(W, T, U, D)$.
  \item \textbf{Identify} the Markov blanket for a given node in a BN. Justify your answer using the definition.
  \item \textbf{Explain} why the Markov blanket includes children's parents but not parents' children.
  \item \textbf{Compute} $P(\text{burglary} \mid \text{alarm})$ using the alarm-network CPTs and Bayes' rule. (Marginalise over earthquake.)
  \item \textbf{Distinguish} discrete and continuous random variables with examples. Which kind do Bayesian networks (in this lecture) use?
  \item \textbf{Describe} the Bernoulli and Binomial distributions. What scenario does each model?
  \item \textbf{Compute} the expected value of a 3-coin-flip experiment scored by ``$\pounds 1$ per heads.''
  \item \textbf{Compare} two random variables defined on the same experiment that have different expected values. Explain why the experiment alone doesn't determine the expectation.
  \item \textbf{Apply} naive Bayes to classify a sentence as spam or ham, given priors and per-word likelihoods of your choosing.
  \item \textbf{State} linearity of expectation. Does it require independence? Why not?
\end{enumerate}

\end{document}
